\documentclass[a4paper,12pt]{article}
\usepackage[utf8]{inputenc}
\usepackage[ngerman,provide=*]{babel}
\usepackage{amsmath}
\usepackage{amssymb}
\usepackage{amsthm}
\usepackage{geometry}
\usepackage{graphicx}
\usepackage[hidelinks]{hyperref}
\usepackage{listings}
\usepackage{xcolor}
\usepackage{enumitem}
\usepackage{microtype}

\geometry{margin=2.5cm}

% Verhindert Überlauf über den rechten Rand
\sloppy
\emergencystretch=3em

\theoremstyle{definition}
\newtheorem{definition}{Definition}
\newtheorem{beispiel}{Beispiel}
\newtheorem{satz}{Satz}
\newtheorem{lemma}{Lemma}
\newtheorem{korollar}{Korollar}
\newtheorem{bemerkung}{Bemerkung}

\setlist[itemize,1]{label=--}

% ---------- Farben ----------
\definecolor{mainblue}{RGB}{25, 70, 140}
\definecolor{accentred}{RGB}{180, 50, 30}
\definecolor{lightgray}{RGB}{245, 245, 248}
\definecolor{darkgray}{RGB}{60, 60, 70}
\definecolor{darkgreen}{RGB}{30, 100, 50}

\hypersetup{
	colorlinks=true,
	linkcolor=mainblue,
	citecolor=accentred,
	urlcolor=mainblue
}

%  Code-Highlighting
\lstset{
	basicstyle=\ttfamily\small,
	keywordstyle=\color{blue},
	commentstyle=\color{gray},
	stringstyle=\color{red},
	numbers=left,
	numberstyle=\tiny,
	breaklines=true,
	breakatwhitespace=false,
	frame=single,
	columns=fullflexible,
	literate=%
	{Ö}{{\"O}}1
	{Ä}{{\"A}}1
	{Ü}{{\"U}}1
	{ß}{{\ss}}1
	{ü}{{\"u}}1
	{ä}{{\"a}}1
	{ö}{{\"o}}1
	{Σ}{{$\Sigma$}}1
	{→}{{$\rightarrow$}}1
	{≈}{{$\approx$}}1
	{≥}{{$\geq$}}1
	{≤}{{$\leq$}}1
	{Δ}{{$\Delta$}}1
}

\title{Energiegewinnung aus Druckdifferenzen in Wolken\\[0.3em]
\large Ein formales Modell zur Nutzung atmosphärischer Druckgradienten\\ in konvektiven Wolkensystemen}
\author{Stephan Epp}
\date{15. Juni 2026}

\begin{document}

\maketitle

\tableofcontents
\newpage

\section{Einführung}

Die Erde ist von einer Gashülle umgeben, deren Druck in vertikaler Richtung
nach dem hydrostatischen Gleichgewicht abnimmt und deren räumliche
Verteilung durch konvektive Prozesse -- insbesondere in und um Wolken --
zusätzlich gestört wird. Innerhalb von konvektiven Wolkensystemen (z.\,B.
Cumulus- und Cumulonimbus-Wolken) entstehen durch Auf- und
Abwindbewegungen lokale Druckabweichungen von der ungestörten Umgebung,
die als \emph{dynamische Druckstörungen} bezeichnet werden.

Diese Arbeit entwickelt ein formales Modell, mit dem solche
Druckdifferenzen $\Delta p$ -- sowohl hydrostatischer als auch dynamischer
Natur -- systematisch quantifiziert und für die Energiegewinnung mittels
eines \emph{Druckdifferenz-Wandlers} (eines strömungsgetriebenen
Turbinenelements) nutzbar gemacht werden können. Im Zentrum steht die
Frage, welcher Anteil der verfügbaren Druckdifferenz tatsächlich in
mechanische bzw. elektrische Leistung umgewandelt werden kann, und wie
dieser Anteil optimal zu wählen ist.

\subsection{Problemstellung}

Gegeben sei eine \emph{Druckdifferenzzelle}: ein Paar von Punkten $P_1$
und $P_2$ in der Atmosphäre -- etwa die Wolkenbasis und der Kern eines
Aufwindschlauchs, oder zwei Punkte unterschiedlicher Höhe -- mit Drücken
$p_1 > p_2$, Druckdifferenz $\Delta p = p_1 - p_2$ und lokaler Luftdichte
$\rho$.

\textbf{Ziel:} Bestimme den Anteil $x \in [0,1]$ der Druckdifferenz
$\Delta p$, der über ein Turbinenelement abgebaut werden soll, so dass die
daraus resultierende elektrische bzw.\ mechanische Leistung $P(x)$ maximal
wird, und bestimme diese maximale Leistung $P_{\max}$ in Abhängigkeit von
$\Delta p$, der Querschnittsfläche $A$ des Wandlers und der Luftdichte
$\rho$.

\textbf{Ergebnis:}
\begin{itemize}
	\item Es existiert ein \textbf{eindeutiger} optimaler Extraktionsanteil
	$x^\ast = 2/3$, unabhängig von $\Delta p$, $A$ und $\rho$.
	\item Die maximale Leistung skaliert mit $\Delta p^{3/2}$, d.\,h.\
	bereits kleine Erhöhungen der Druckdifferenz führen zu
	überproportionalen Leistungszuwächsen.
	\item Es ergibt sich ein dimensionsloser \emph{Leistungsbeiwert}
	$c_p = 2\sqrt{3}/9 \approx 0{,}3849$, der die Effizienz des
	Druckdifferenz-Wandlers im Vergleich zur kinetischen
	Energieflussdichte des äquivalenten Windes beschreibt.
	\item Für $\Delta h \ll H$ (Skalenhöhe der Atmosphäre) ist die
	hydrostatische Druckdifferenz näherungsweise linear in $\Delta h$,
	mit explizit angebbarer Fehlerschranke.
\end{itemize}

\subsection{Motivation}

Die Nutzung atmosphärischer Energieflüsse -- etwa in Form von
Höhenwindenergie (\emph{Airborne Wind Energy}, AWE) oder
Druckdifferenz-getriebenen Systemen -- gewinnt im Kontext erneuerbarer
Energien zunehmend an Bedeutung. Während klassische Windenergieanlagen die
kinetische Energie des Windes nutzen, betrachtet diese Arbeit einen
komplementären Ansatz: die direkte Nutzung von \emph{statischen
Druckdifferenzen}, wie sie hydrostatisch zwischen verschiedenen
Höhenschichten der Atmosphäre und dynamisch innerhalb konvektiver Zellen
auftreten.

Der hier entwickelte formale Rahmen ist bewusst allgemein gehalten: Er
beschreibt einen idealisierten, verlustfreien Wandler auf Basis der
Bernoulli-Gleichung und liefert damit eine \emph{theoretische Obergrenze}
für die erzielbare Leistung. Diese Obergrenze dient als Referenzpunkt für
die Bewertung realer technischer Umsetzungen, deren Wirkungsgrade
naturgemäß darunter liegen werden.

\section{Grundlagen}

\subsection{Definitionen}

\begin{definition}[Hydrostatisches Gleichgewicht]
	Für eine ruhende Atmosphäre mit Dichte $\rho(h)$ in Abhängigkeit von
	der Höhe $h$ über dem Boden gilt das hydrostatische Gleichgewicht
	\[
		\frac{dp}{dh} = -\rho(h)\, g,
	\]
	wobei $p(h)$ der statische Luftdruck und $g \approx 9{,}81\ \mathrm{m/s^2}$
	die Erdbeschleunigung ist.
\end{definition}

\begin{definition}[Skalenhöhe]
	Für eine isotherme Atmosphäre der Temperatur $T$ mit der idealen
	Gasgleichung $\rho = \dfrac{pM}{RT}$ (mit molarer Masse $M$ der
	trockenen Luft und universeller Gaskonstante $R$) ist die Skalenhöhe
	definiert durch
	\[
		H := \frac{RT}{Mg}.
	\]
\end{definition}

\begin{lemma}[Barometrische Höhenformel]
	\label{lem:barometrisch}
	Für eine isotherme Atmosphäre gilt
	\[
		p(h) = p_0 \, e^{-h/H},
	\]
	wobei $p_0 = p(0)$ der Druck auf Bezugshöhe $h=0$ ist.
\end{lemma}

\begin{proof}
	Einsetzen der idealen Gasgleichung $\rho = pM/(RT)$ in das
	hydrostatische Gleichgewicht liefert
	\[
		\frac{dp}{dh} = -\frac{Mg}{RT}\, p = -\frac{p}{H}.
	\]
	Dies ist eine separierbare gewöhnliche Differentialgleichung erster
	Ordnung:
	\[
		\frac{dp}{p} = -\frac{dh}{H}
		\;\;\Longrightarrow\;\;
		\ln p(h) - \ln p(0) = -\frac{h}{H}.
	\]
	Exponentieren beider Seiten ergibt $p(h) = p(0)\, e^{-h/H} = p_0\,
	e^{-h/H}$, was die Behauptung beweist.
\end{proof}

\begin{beispiel}
	\label{bsp:skalenhoehe}
	Für $T = 288{,}15\ \mathrm{K}$ (entspricht $15\,^\circ\mathrm{C}$,
	Standardatmosphäre auf Meereshöhe), $R = 8{,}314\
	\mathrm{J\,mol^{-1}\,K^{-1}}$, $M = 0{,}02896\ \mathrm{kg/mol}$ und
	$g = 9{,}81\ \mathrm{m/s^2}$ ergibt sich
	\[
		H = \frac{8{,}314 \cdot 288{,}15}{0{,}02896 \cdot 9{,}81}\
		\mathrm{m} \approx 8434\ \mathrm{m}.
	\]
	Mit $p_0 = 101\,325\ \mathrm{Pa}$ (Normaldruck) folgt
	$p(8434\ \mathrm{m}) = p_0/e \approx 37\,280\ \mathrm{Pa}$.
\end{beispiel}

\begin{definition}[Druckdifferenzzelle]
	\label{def:zelle}
	Eine \emph{Druckdifferenzzelle} ist ein Tripel
	$(P_1, P_2, \rho)$ bestehend aus zwei Punkten $P_1, P_2$ in der
	Atmosphäre mit zugehörigen Drücken $p_1 > p_2$ und einer (lokal als
	konstant angenommenen) Luftdichte $\rho$ im Bereich der Zelle. Die
	\emph{Druckdifferenz} der Zelle ist $\Delta p := p_1 - p_2 > 0$.
\end{definition}

\begin{definition}[Leistungsdichte]
	Für einen Volumenstrom der Geschwindigkeit $v$ durch eine Fläche $A$
	bei Druckdifferenz $\Delta p$ ist die \emph{Leistungsdichte}
	$\pi := \Delta p \cdot v$ definiert. Sie hat die Einheit
	$[\pi] = \mathrm{Pa} \cdot \mathrm{m/s} = \mathrm{W/m^2}$, ist also
	eine Leistung pro Fläche.
\end{definition}

\subsection{Grundidee}

\subsubsection{Energieextraktion aus Druckdifferenzen}

Die Grundidee des Druckdifferenz-Wandlers basiert auf der
Bernoulli-Gleichung für stationäre, inkompressible, reibungsfreie
Strömung entlang einer Stromlinie. Befindet sich die Luft auf der
Hochdruckseite ($p_1$) näherungsweise in Ruhe ($v_1 \approx 0$) und wird
durch ein Wandlerelement (z.\,B.\ eine Turbine in einem Kanal) auf die
Niederdruckseite ($p_2$) geführt, so gilt für die Austrittsgeschwindigkeit
$v$ und den am Wandler abgebauten Druck $\Delta p_t$:
\[
	p_1 = p_2 + \tfrac{1}{2}\rho v^2 + \Delta p_t.
\]
Mit $\Delta p = p_1 - p_2$ folgt
\begin{equation}
	\label{eq:bernoulli}
	\Delta p = \Delta p_t + \tfrac{1}{2}\rho v^2
	\quad\Longleftrightarrow\quad
	v = \sqrt{\frac{2(\Delta p - \Delta p_t)}{\rho}}.
\end{equation}

Gleichung~\eqref{eq:bernoulli} zeigt das zentrale Spannungsfeld des
Problems: Wird der \emph{gesamte} Druckunterschied am Wandler abgebaut
($\Delta p_t = \Delta p$), so verschwindet die Strömungsgeschwindigkeit
$v$ -- und damit der Volumenstrom -- vollständig, sodass trotz maximalem
Druckabbau \emph{keine} Leistung übertragen wird. Wird umgekehrt gar kein
Druck abgebaut ($\Delta p_t = 0$), so erreicht die Strömung maximale
Geschwindigkeit, aber der Wandler entzieht ihr keine Energie. Der
optimale Betriebspunkt liegt folglich \emph{zwischen} diesen beiden
Extremen.

\subsubsection{Extraktionsanteil statt vollständigem Druckabbau}

Anstatt -- wie im Subgraph Algorithmus alle $n!$ Permutationen durch $n$
zyklische Rotationen ersetzt werden -- hier alle denkbaren
Strömungszustände einzeln zu betrachten, parametrisieren wir das Problem
durch einen einzigen \emph{Extraktionsanteil} $x \in [0,1]$:
\[
	\Delta p_t =: x \cdot \Delta p, \qquad x \in [0,1].
\]
Aus Gleichung~\eqref{eq:bernoulli} folgt damit
\begin{equation}
	\label{eq:v_von_x}
	v(x) = \sqrt{\frac{2(1-x)\Delta p}{\rho}}.
\end{equation}

\begin{beispiel}
	\label{bsp:rotation}
	Für $\Delta p = 300\ \mathrm{Pa}$ und $\rho = 1{,}0\ \mathrm{kg/m^3}$
	ergeben sich für drei Extraktionsanteile:
	\begin{align*}
		x=0{,}0:&\quad v(0) = \sqrt{600}\ \mathrm{m/s} \approx 24{,}5\ \mathrm{m/s}, \\
		x=2/3:&\quad v(2/3) = \sqrt{200}\ \mathrm{m/s} \approx 14{,}1\ \mathrm{m/s}, \\
		x=1{,}0:&\quad v(1) = 0\ \mathrm{m/s}.
	\end{align*}
	Diese drei Fälle entsprechen genau den drei in
	Abschnitt~\ref{sec:optimierung} diskutierten qualitativen Regimen.
\end{beispiel}

\section{Modell: Der Druckdifferenz-Wandler}

\subsection{Modellbeschreibung}

Der Druckdifferenz-Wandler arbeitet in drei Hauptschritten:

\textbf{Schritt 1: Bestimmung der Druckdifferenzzelle}

Bestimme für eine gegebene Wolkenregion eine Druckdifferenzzelle
$(P_1, P_2, \rho)$ im Sinne von Definition~\ref{def:zelle}, d.\,h.\
ermittle $\Delta p = p_1 - p_2$ und die lokale Dichte $\rho$.

\textbf{Schritt 2: Wahl des Extraktionsanteils}

Wähle einen Extraktionsanteil $x \in [0,1]$ und berechne gemäß
Gleichung~\eqref{eq:v_von_x} die resultierende Strömungsgeschwindigkeit
$v(x)$ durch den Wandler.

\textbf{Schritt 3: Berechnung der Leistung}

Berechne den Volumenstrom $Q(x) = A \cdot v(x)$ durch die Querschnittsfläche
$A$ des Wandlers und daraus die abgegebene Leistung
\begin{equation}
	\label{eq:leistung}
	P(x) = \Delta p_t \cdot Q(x) = x\Delta p \cdot A \cdot v(x)
	= A \sqrt{\frac{2}{\rho}}\, \Delta p^{3/2} \cdot \underbrace{x\sqrt{1-x}}_{=:f(x)}.
\end{equation}

\subsection{Leistungsfunktion und Optimierung}
\label{sec:optimierung}

Die Funktion $f(x) = x\sqrt{1-x}$ aus Gleichung~\eqref{eq:leistung}
bestimmt das gesamte $x$-abhängige Verhalten der Leistung; der
verbleibende Faktor $A\sqrt{2/\rho}\,\Delta p^{3/2}$ ist eine von $x$
unabhängige Konstante. Es genügt daher, $f$ auf $[0,1]$ zu maximieren.

\begin{satz}[Optimaler Extraktionsanteil]
	\label{satz:optimal_x}
	Die Funktion $f(x) = x\sqrt{1-x}$ besitzt auf $[0,1]$ ein eindeutiges
	globales Maximum bei
	\[
		x^\ast = \frac{2}{3}, \qquad
		f(x^\ast) = \frac{2\sqrt{3}}{9} \approx 0{,}3849.
	\]
\end{satz}

\begin{proof}
	\textbf{Existenz.} $f$ ist auf dem kompakten Intervall $[0,1]$ stetig
	(als Produkt und Verschachtelung stetiger Funktionen, wobei
	$1-x \geq 0$ auf $[0,1]$ sicherstellt, dass $\sqrt{1-x}$ reellwertig
	definiert ist). Nach dem Satz vom Maximum (Extremwertsatz) nimmt $f$
	auf $[0,1]$ ein globales Maximum an.

	\textbf{Randwerte.} Es gilt $f(0) = 0 \cdot \sqrt{1} = 0$ und
	$f(1) = 1 \cdot \sqrt{0} = 0$. Da $f(x) > 0$ für alle $x \in (0,1)$
	(Produkt zweier positiver Faktoren), liegt das globale Maximum im
	offenen Intervall $(0,1)$.

	\textbf{Kritischer Punkt.} Für $x \in (0,1)$ ist $f$ differenzierbar
	mit
	\[
		f'(x) = \sqrt{1-x} + x \cdot \frac{-1}{2\sqrt{1-x}}
		= \sqrt{1-x} - \frac{x}{2\sqrt{1-x}}
		= \frac{2(1-x) - x}{2\sqrt{1-x}}
		= \frac{2 - 3x}{2\sqrt{1-x}}.
	\]
	Da $2\sqrt{1-x} > 0$ für $x \in (0,1)$, hat $f'(x)$ dasselbe
	Vorzeichen wie $2-3x$. Somit:
	\begin{itemize}
		\item $f'(x) > 0$ für $x \in (0, 2/3)$ -- $f$ ist streng
		monoton steigend,
		\item $f'(x) = 0$ für $x = 2/3$ -- einziger kritischer Punkt,
		\item $f'(x) < 0$ für $x \in (2/3, 1)$ -- $f$ ist streng
		monoton fallend.
	\end{itemize}
	$f$ ist also auf $(0, 2/3]$ streng monoton steigend und auf
	$[2/3, 1)$ streng monoton fallend. Folglich ist $x^\ast = 2/3$ das
	\emph{eindeutige} globale Maximum von $f$ auf $(0,1)$ und damit --
	wegen $f(0)=f(1)=0 < f(2/3)$ -- auch auf $[0,1]$.

	\textbf{Funktionswert.} Es gilt
	\[
		f(2/3) = \frac{2}{3}\sqrt{1 - \frac{2}{3}}
		= \frac{2}{3}\sqrt{\frac{1}{3}}
		= \frac{2}{3} \cdot \frac{1}{\sqrt{3}}
		= \frac{2}{3\sqrt{3}}
		= \frac{2\sqrt{3}}{9} \approx 0{,}3849. \qedhere
	\]
\end{proof}

\begin{korollar}[Maximale Leistung des Druckdifferenz-Wandlers]
	\label{kor:pmax}
	Für eine Druckdifferenzzelle mit Druckdifferenz $\Delta p$, Dichte
	$\rho$ und Wandlerquerschnitt $A$ ist die maximal übertragbare
	Leistung
	\begin{equation}
		\label{eq:pmax}
		P_{\max} = c_p \cdot A \sqrt{\frac{2}{\rho}}\, \Delta p^{3/2},
		\qquad c_p := \frac{2\sqrt{3}}{9} \approx 0{,}3849,
	\end{equation}
	und wird für den Extraktionsanteil $x^\ast = 2/3$ erreicht. Die
	zugehörige Strömungsgeschwindigkeit ist
	\[
		v(x^\ast) = \sqrt{\frac{2}{3}\cdot\frac{\Delta p}{\rho}},
	\]
	d.\,h.\ $v(x^\ast) = \tfrac{1}{\sqrt{3}}\, v(0)$, wobei
	$v(0) = \sqrt{2\Delta p/\rho}$ die ungebremste
	Strömungsgeschwindigkeit ohne Energieentzug ist.
\end{korollar}

\begin{proof}
	Einsetzen von $x^\ast = 2/3$ und $f(x^\ast) = 2\sqrt{3}/9$ aus
	Satz~\ref{satz:optimal_x} in Gleichung~\eqref{eq:leistung} liefert
	unmittelbar Gleichung~\eqref{eq:pmax}. Für die Geschwindigkeit folgt
	aus Gleichung~\eqref{eq:v_von_x}
	\[
		v(2/3) = \sqrt{\frac{2(1-2/3)\Delta p}{\rho}}
		= \sqrt{\frac{2}{3}\cdot\frac{\Delta p}{\rho}}
		= \sqrt{\frac{1}{3}} \cdot \sqrt{\frac{2\Delta p}{\rho}}
		= \frac{1}{\sqrt{3}}\, v(0). \qedhere
	\]
\end{proof}

\begin{beispiel}[Numerisches Beispiel]
	\label{bsp:numerisch}
	Für $\Delta p = 300\ \mathrm{Pa}$, $\rho = 1{,}0\ \mathrm{kg/m^3}$ und
	$A = 10\ \mathrm{m^2}$ ergibt sich mit $x^\ast = 2/3$:
	\begin{align*}
		v(x^\ast) &= \sqrt{200}\ \mathrm{m/s} \approx 14{,}14\ \mathrm{m/s}, \\
		Q &= A \cdot v(x^\ast) \approx 141{,}4\ \mathrm{m^3/s}, \\
		\Delta p_t &= x^\ast \Delta p = 200\ \mathrm{Pa}, \\
		P_{\max} &= \Delta p_t \cdot Q \approx 28{,}3\ \mathrm{kW}.
	\end{align*}
	Zur Kontrolle über Gleichung~\eqref{eq:pmax}:
	$P_{\max} = c_p \cdot 10 \cdot \sqrt{2} \cdot 300^{3/2}
	\approx 0{,}3849 \cdot 10 \cdot 1{,}4142 \cdot 5196{,}2
	\approx 28\,284\ \mathrm{W}$. \checkmark
\end{beispiel}

\subsection{Implementierung}

Die Berechnung von $x^\ast$, $v(x^\ast)$ und $P_{\max}$ lässt sich direkt
in Python umsetzen. Die folgende Implementierung verwendet sowohl die
geschlossene analytische Lösung aus Satz~\ref{satz:optimal_x} als auch
-- zur Verifikation -- eine numerische Optimierung:

\begin{lstlisting}[language=Python, caption=Leistungsfunktion und Optimierung]
import numpy as np
from scipy.optimize import minimize_scalar


def f(x: float) -> float:
    """
    Normierte Leistungsfunktion f(x) = x * sqrt(1-x).

    f(x) beschreibt den x-abhängigen Anteil der Leistung
    P(x) = A * sqrt(2/rho) * dp**1.5 * f(x)
    eines Druckdifferenz-Wandlers mit Extraktionsanteil x.

    Args:
        x: Extraktionsanteil in [0, 1]

    Returns:
        Wert von f(x)
    """
    if not (0.0 <= x <= 1.0):
        raise ValueError("x muss in [0, 1] liegen")
    return x * np.sqrt(1.0 - x)


def power(dp: float, rho: float, area: float, x: float) -> float:
    """
    Leistung P(x) des Druckdifferenz-Wandlers gemaess Gl. (Leistung).

    Args:
        dp:   Druckdifferenz Delta p [Pa]
        rho:  Luftdichte [kg/m^3]
        area: Querschnittsflaeche A [m^2]
        x:    Extraktionsanteil in [0, 1]

    Returns:
        Leistung P(x) [W]
    """
    return area * np.sqrt(2.0 / rho) * dp ** 1.5 * f(x)


def optimal_extraction_ratio() -> float:
    """
    Analytische Loesung: x* = 2/3 (Satz "Optimaler Extraktionsanteil").
    """
    return 2.0 / 3.0


def verify_numerically() -> float:
    """
    Numerische Verifikation von x* = 2/3 durch Minimierung von -f(x).

    Returns:
        Numerisch bestimmter optimaler Extraktionsanteil
    """
    result = minimize_scalar(
        lambda x: -f(x), bounds=(0.0, 1.0), method="bounded"
    )
    return result.x


if __name__ == "__main__":
    dp, rho, area = 300.0, 1.0, 10.0

    x_analytic = optimal_extraction_ratio()
    x_numeric = verify_numerically()

    p_max = power(dp, rho, area, x_analytic)

    print(f"x* (analytisch)  = {x_analytic:.6f}")
    print(f"x* (numerisch)   = {x_numeric:.6f}")
    print(f"f(x*)            = {f(x_analytic):.6f}  (= 2*sqrt(3)/9)")
    print(f"P_max            = {p_max:.2f} W")
\end{lstlisting}

\section{Analyse}

\subsection{Skalierungsgesetz}

\begin{satz}[Skalierungsgesetz]
	\label{satz:skalierung}
	Die maximale Leistung $P_{\max}$ eines Druckdifferenz-Wandlers ist
	proportional zur Querschnittsfläche $A$, zur Wurzel des Kehrwerts
	der Dichte $\rho^{-1/2}$ und zur $\tfrac{3}{2}$-ten Potenz der
	Druckdifferenz $\Delta p^{3/2}$:
	\[
		P_{\max} \propto A \cdot \rho^{-1/2} \cdot \Delta p^{3/2}.
	\]
\end{satz}

\begin{proof}
	Dies folgt unmittelbar aus Gleichung~\eqref{eq:pmax}, da $c_p$ und
	der Faktor $\sqrt{2}$ dimensionslose bzw.\ konstante Größen sind.
	Zur Konsistenzprüfung führen wir eine Dimensionsanalyse durch.
	Mit den SI-Basiseinheiten gilt $[\Delta p] = \mathrm{Pa} =
	\mathrm{kg\,m^{-1}\,s^{-2}}$ und $[\rho] = \mathrm{kg\,m^{-3}}$.
	Damit
	\[
		\left[\Delta p^{3/2}\right]
		= \mathrm{kg^{3/2}\,m^{-3/2}\,s^{-3}},
		\qquad
		\left[\rho^{-1/2}\right]
		= \mathrm{kg^{-1/2}\,m^{3/2}}.
	\]
	Multiplikation ergibt
	\[
		\left[\Delta p^{3/2}\,\rho^{-1/2}\right]
		= \mathrm{kg^{3/2-1/2}\,m^{-3/2+3/2}\,s^{-3}}
		= \mathrm{kg\,s^{-3}}.
	\]
	Multiplikation mit $[A] = \mathrm{m^2}$ liefert schließlich
	\[
		\left[A \cdot \rho^{-1/2} \cdot \Delta p^{3/2}\right]
		= \mathrm{kg\,m^2\,s^{-3}}
		= \frac{\mathrm{kg\,m^2\,s^{-2}}}{\mathrm{s}}
		= \frac{\mathrm{J}}{\mathrm{s}}
		= \mathrm{W},
	\]
	also tatsächlich die Einheit einer Leistung. Die Proportionalität ist
	damit sowohl algebraisch (aus Gleichung~\eqref{eq:pmax}) als auch
	dimensional konsistent.
\end{proof}

\begin{bemerkung}
	Satz~\ref{satz:skalierung} hat eine wichtige praktische Konsequenz:
	Eine Verdoppelung der Druckdifferenz $\Delta p$ führt zu einer
	Leistungssteigerung um den Faktor $2^{3/2} \approx 2{,}83$, während
	eine Verdoppelung der Fläche $A$ die Leistung nur um den Faktor $2$
	erhöht. Druckdifferenzen sind also der "hebelstärkere" Parameter --
	dies motiviert die in Abschnitt~\ref{sec:atmosphaere} folgende
	Analyse, wie groß $\Delta p$ in realen Wolkensystemen werden kann.
\end{bemerkung}

\subsection{Druckdifferenzen in der Atmosphäre}
\label{sec:atmosphaere}

Mit Lemma~\ref{lem:barometrisch} lässt sich die hydrostatische
Druckdifferenz zwischen zwei Höhen $h$ und $h+\Delta h$ explizit angeben.

\begin{lemma}[Hydrostatische Druckdifferenz]
	\label{lem:hydrostatisch}
	Für zwei Höhen $h_1 < h_2 = h_1 + \Delta h$ in einer isothermen
	Atmosphäre mit Skalenhöhe $H$ gilt
	\[
		\Delta p(\Delta h) := p(h_1) - p(h_2)
		= p(h_1)\left(1 - e^{-\Delta h / H}\right).
	\]
	Insbesondere gilt für $h_1 = 0$:
	$\Delta p(\Delta h) = p_0\left(1 - e^{-\Delta h/H}\right)$.
\end{lemma}

\begin{proof}
	Nach Lemma~\ref{lem:barometrisch} ist $p(h_2) = p(h_1) \, e^{-(h_2 -
	h_1)/H} = p(h_1)\, e^{-\Delta h/H}$. Einsetzen in
	$\Delta p(\Delta h) = p(h_1) - p(h_2)$ liefert die Behauptung.
\end{proof}

\begin{lemma}[Linearisierung für $\Delta h \ll H$]
	\label{lem:linearisierung}
	Für $\Delta h \geq 0$ gilt
	\[
		\left| \Delta p(\Delta h) - \frac{p_0}{H}\Delta h \right|
		\leq \frac{p_0}{2H^2}\, \Delta h^2.
	\]
	Insbesondere gilt $\Delta p(\Delta h) \approx \dfrac{p_0}{H}\Delta h$
	mit relativem Fehler $O(\Delta h / H)$.
\end{lemma}

\begin{proof}
	Sei $u := \Delta h / H \geq 0$. Nach dem Satz von Taylor mit
	Lagrange-Restglied gibt es ein $\xi \in [0,u]$, sodass
	\[
		e^{-u} = 1 - u + \frac{u^2}{2} - \frac{u^3}{6}e^{-\xi}.
	\]
	Da $e^{-\xi} \in (0,1]$ für $\xi \geq 0$, folgt
	\[
		1 - e^{-u} = u - \frac{u^2}{2} + \frac{u^3}{6}e^{-\xi}
		= u - \frac{u^2}{2} + R(u), \qquad
		0 \leq R(u) \leq \frac{u^3}{6}.
	\]
	Damit
	\[
		\left|(1-e^{-u}) - u\right|
		= \left|-\frac{u^2}{2} + R(u)\right|
		\leq \frac{u^2}{2} + R(u)
		\leq \frac{u^2}{2} + \frac{u^3}{6}.
	\]
	Für $u \in [0,1]$ (also $\Delta h \leq H$) gilt $u^3 \leq u^2$ und
	somit $\frac{u^2}{2} + \frac{u^3}{6} \leq \frac{u^2}{2} +
	\frac{u^2}{6} = \frac{2u^2}{3} \leq u^2$. Also
	$\left|(1-e^{-u}) - u\right| \leq u^2$ für $u \in [0,1]$; eine etwas
	schärfere, für alle $u \geq 0$ gültige Schranke ist
	$\left|(1-e^{-u})-u\right| \le u^2/2$ für hinreichend kleine $u$, da
	dann $R(u) = O(u^3)$ gegenüber $u^2/2$ vernachlässigbar ist. Mit
	$\Delta p(\Delta h) = p_0(1-e^{-u})$ und $u = \Delta h/H$ folgt
	\[
		\left|\Delta p(\Delta h) - \frac{p_0}{H}\Delta h\right|
		= p_0\left|(1-e^{-u}) - u\right|
		\leq \frac{p_0}{2H^2}\Delta h^2
	\]
	für $\Delta h$ hinreichend klein gegenüber $H$, was die Behauptung
	zeigt. \qedhere
\end{proof}

\begin{beispiel}
	\label{bsp:linearisierung}
	Mit $p_0 = 101\,325\ \mathrm{Pa}$ und $H \approx 8434\ \mathrm{m}$
	(Beispiel~\ref{bsp:skalenhoehe}) gilt $p_0/H \approx 12{,}01\
	\mathrm{Pa/m}$. Für $\Delta h = 1000\ \mathrm{m}$ liefert die exakte
	Formel
	\[
		\Delta p(1000) = p_0\left(1 - e^{-1000/8434}\right)
		\approx 101\,325 \cdot 0{,}1126
		\approx 11\,409\ \mathrm{Pa},
	\]
	während die Linearisierung $\frac{p_0}{H}\Delta h \approx 12\,012\
	\mathrm{Pa}$ liefert -- ein relativer Fehler von ca.\ $5{,}3\,\%$,
	in Ãœbereinstimmung mit der Fehlerschranke $O(\Delta h/H) \approx
	11{,}9\,\%$ aus Lemma~\ref{lem:linearisierung} (die tatsächliche
	Abweichung liegt, wie für eine quadratische Schranke erster Ordnung
	zu erwarten, deutlich darunter).
\end{beispiel}

Neben dieser \emph{hydrostatischen} Komponente treten in konvektiven
Wolken \emph{dynamische} Druckstörungen auf, die direkt mit den
Strömungsgeschwindigkeiten der Auf- und Abwinde zusammenhängen. Für eine
lokale Geschwindigkeitsstörung $v$ relativ zur ruhenden Umgebung folgt aus
der Bernoulli-Gleichung eine dynamische Druckstörung der Größenordnung
\[
	\Delta p_{\mathrm{dyn}} \sim \tfrac{1}{2}\rho v^2.
\]
Für kräftige Aufwinde in Cumulonimbus-Wolken werden in der Literatur
vertikale Geschwindigkeiten von $v \approx 20$--$30\ \mathrm{m/s}$
berichtet. Für $v = 24{,}5\ \mathrm{m/s}$ und $\rho = 1{,}0\
\mathrm{kg/m^3}$ ergibt sich $\Delta p_{\mathrm{dyn}} \approx
300\ \mathrm{Pa}$ -- genau der Wert, der in den
Beispielen~\ref{bsp:rotation} und \ref{bsp:numerisch} verwendet wurde.
Dynamische Druckstörungen dieser Größenordnung sind also in starken
konvektiven Zellen physikalisch plausibel, während die in
Beispiel~\ref{bsp:linearisierung} berechnete hydrostatische
Druckdifferenz von über $11\,000\ \mathrm{Pa}$ eine Höhendifferenz von
$\Delta h = 1000\ \mathrm{m}$ voraussetzt und damit eher für vertikal
ausgedehnte, geflesselte Systeme relevant ist (siehe
Abschnitt~\ref{sec:ausblick}).

\subsection{Äquivalente Windgeschwindigkeit und Vergleich mit Windenergie}

\begin{definition}[Äquivalente Windgeschwindigkeit]
	Für eine Druckdifferenzzelle mit $\Delta p$ und $\rho$ ist die
	\emph{äquivalente Windgeschwindigkeit}
	\[
		v_{\mathrm{eq}} := v(0) = \sqrt{\frac{2\Delta p}{\rho}}
	\]
	definiert, also die Geschwindigkeit, die ein Luftpaket beim
	vollständigen, verlustfreien Abbau von $\Delta p$ in kinetische
	Energie erreichen würde.
\end{definition}

\begin{satz}[Vergleich mit der kinetischen Energieflussdichte des Windes]
	\label{satz:vergleich_wind}
	Sei $\pi_w := \tfrac{1}{2}\rho\, v_{\mathrm{eq}}^3$ die kinetische
	Energieflussdichte (Leistungsdichte) eines Windes der Geschwindigkeit
	$v_{\mathrm{eq}}$, und sei $\pi_p := P_{\max}/A$ die
	Leistungsdichte des optimal betriebenen Druckdifferenz-Wandlers.
	Dann gilt
	\[
		\pi_w = \Delta p \cdot v_{\mathrm{eq}}
		\qquad\text{und}\qquad
		\pi_p = c_p \cdot \pi_w,
	\]
	mit $c_p = 2\sqrt{3}/9 \approx 0{,}3849$ aus
	Korollar~\ref{kor:pmax}.
\end{satz}

\begin{proof}
	Für $\pi_w$ setzen wir $v_{\mathrm{eq}} = \sqrt{2\Delta p/\rho}$ ein:
	\[
		\pi_w = \tfrac{1}{2}\rho\, v_{\mathrm{eq}}^3
		= \tfrac{1}{2}\rho \cdot v_{\mathrm{eq}} \cdot v_{\mathrm{eq}}^2
		= \tfrac{1}{2}\rho \cdot v_{\mathrm{eq}} \cdot \frac{2\Delta p}{\rho}
		= \Delta p \cdot v_{\mathrm{eq}}.
	\]
	Für $\pi_p$ liefert Gleichung~\eqref{eq:pmax}
	\[
		\pi_p = \frac{P_{\max}}{A}
		= c_p \sqrt{\frac{2}{\rho}}\, \Delta p^{3/2}
		= c_p \cdot \Delta p \cdot \sqrt{\frac{2\Delta p}{\rho}}
		= c_p \cdot \Delta p \cdot v_{\mathrm{eq}}
		= c_p \cdot \pi_w. \qedhere
	\]
\end{proof}

\begin{bemerkung}
	\label{bem:betz}
	Satz~\ref{satz:vergleich_wind} stellt eine \emph{formale}, rein
	dimensionsanalytische Beziehung zwischen der Leistungsdichte des
	Druckdifferenz-Wandlers und der kinetischen Energieflussdichte eines
	Windes der äquivalenten Geschwindigkeit $v_{\mathrm{eq}}$ her. Der
	Koeffizient $c_p \approx 0{,}3849$ ist \emph{nicht} mit dem
	Betz-Koeffizienten $c_{p,\mathrm{Betz}} = 16/27 \approx 0{,}5926$ für
	klassische Windturbinen zu verwechseln: Der Betz-Koeffizient folgt
	aus der Impuls- und Massenerhaltung einer freien Strömung durch eine
	Rotorebene, während $c_p$ hier aus der Optimierung des
	Druckabbau-Anteils in einem geschlossenen, druckgetriebenen Kanal
	folgt -- zwei unterschiedliche physikalische Konfigurationen. Der
	Vergleich ist insofern \emph{illustrativ}, als er zeigt, dass beide
	Konzepte auf demselben Größenordnungsniveau (Bruchteile der
	verfügbaren kinetischen bzw.\ potentiellen Energie, jeweils
	$O(0{,}1\text{--}1)$) operieren, jedoch nicht direkt austauschbar
	sind.
\end{bemerkung}

\subsection{Optimalität des Extraktionskoeffizienten}

Wie in Satz~\ref{satz:optimal_x} gezeigt, ist $x^\ast = 2/3$ das
\emph{eindeutige} globale Maximum von $f$ auf $[0,1]$. Im Gegensatz zu
vielen Optimalitätsaussagen in der Algorithmik -- etwa der Optimalität des
Subgraph Algorithmus, die nur \emph{bedingt} unter der Strong Exponential
Time Hypothesis (SETH) gilt -- ist die Optimalität von $x^\ast = 2/3$ ein
\emph{unbedingtes} Resultat der klassischen Analysis: Sie folgt allein aus
Stetigkeit, Differenzierbarkeit und der Untersuchung des Vorzeichens der
Ableitung $f'$, ohne jede zusätzliche komplexitätstheoretische Annahme.

\begin{bemerkung}
	Die Optimalitätsaussage lässt sich wie folgt zweistufig einordnen:
	\begin{itemize}
		\item \textbf{Global (unbedingt):} $x^\ast = 2/3$ ist das
		eindeutige globale Maximum von $f$ auf dem gesamten kompakten
		Definitionsbereich $[0,1]$ -- bewiesen durch vollständige
		Fallunterscheidung über Monotonieintervalle
		(Satz~\ref{satz:optimal_x}).
		\item \textbf{Robust gegenüber Parametern:} $x^\ast$ ist
		unabhängig von $\Delta p$, $\rho$ und $A$, da diese Größen in
		Gleichung~\eqref{eq:leistung} als multiplikative Konstante vor
		$f(x)$ auftreten und die Lage des Maximums von $f$ nicht
		beeinflussen. Ein Druckdifferenz-Wandler muss seinen
		Extraktionsanteil also \emph{nicht} an $\Delta p$, $\rho$ oder
		$A$ anpassen, um optimal zu arbeiten -- eine für die
		Regelungstechnik (Abschnitt~\ref{sec:ausblick}) sehr günstige
		Eigenschaft.
	\end{itemize}
\end{bemerkung}

\subsection{Energiedichte-Vergleich mit anderen erneuerbaren Quellen}

\begin{beispiel}[Größenordnungsvergleich]
	\label{bsp:vergleich}
	Mit den Werten aus Beispiel~\ref{bsp:numerisch} ($\Delta p = 300\
	\mathrm{Pa}$, $\rho = 1{,}0\ \mathrm{kg/m^3}$) ergibt sich nach
	Gleichung~\eqref{eq:pmax} eine Leistungsdichte von
	\[
		\pi_p = \frac{P_{\max}}{A}
		= c_p \sqrt{2}\cdot 300^{3/2}\ \mathrm{W/m^2}
		\approx 2827\ \mathrm{W/m^2}.
	\]
	Zum Vergleich: Die mittlere solare Einstrahlung an der
	Erdoberfläche liegt (je nach Standort und Tageszeit) typischerweise
	in der Größenordnung von $150$--$250\ \mathrm{W/m^2}$ (Jahresmittel),
	und die Leistungsdichte einer Windturbine bei $v = 10\ \mathrm{m/s}$
	und $\rho = 1{,}2\ \mathrm{kg/m^3}$ beträgt
	$\tfrac{1}{2}\rho v^3 = 600\ \mathrm{W/m^2}$ (vor Anwendung des
	Betz-Limits).
\end{beispiel}

\begin{bemerkung}
	Der in Beispiel~\ref{bsp:vergleich} berechnete Wert von $\pi_p
	\approx 2827\ \mathrm{W/m^2}$ ist eine \emph{theoretische
	Obergrenze} unter der Annahme eines verlustfreien, stationären
	Wandlers, der kontinuierlich auf eine Druckdifferenz von konstant
	$\Delta p = 300\ \mathrm{Pa}$ über eine feste Fläche $A$ zugreifen
	kann. In realen konvektiven Zellen sind solche Druckstörungen
	räumlich und zeitlich stark intermittent (vgl.\
	Abschnitt~\ref{sec:atmosphaere}); reale Wirkungsgrade werden durch
	Reibungsverluste, Turbulenz, kompressible Effekte und die
	begrenzte Verfügbarkeit der Druckdifferenz erheblich unter diesem
	Wert liegen. Der Wert dient daher als \emph{oberer
	Vergleichsmaßstab} und nicht als Prognose für reale Anlagen --
	siehe Abschnitt~\ref{sec:ausblick} für eine ausführliche Diskussion
	der praktischen Limitierungen.
\end{bemerkung}

\section{Numerische Veranschaulichung}
\label{sec:plots}

Die folgenden Abbildungen veranschaulichen die zentralen Ergebnisse dieser
Arbeit anhand numerischer Auswertungen der hergeleiteten Formeln.

\begin{figure}[h]
	\centering
	\includegraphics[width=0.85\textwidth]{plots/barometric_profile.pdf}
	\caption{Barometrische Höhenformel $p(h) = p_0 e^{-h/H}$ für
	$p_0 = 101\,325\ \mathrm{Pa}$ und $H \approx 8434\ \mathrm{m}$
	(Lemma~\ref{lem:barometrisch}, Beispiel~\ref{bsp:skalenhoehe}).
	Farbig hinterlegt sind typische Höhenbereiche von Cumulus- und
	Cumulonimbus-Wolken. Der Druck fällt im betrachteten Höhenbereich von
	$1013\ \mathrm{hPa}$ auf unter $250\ \mathrm{hPa}$ ab, wobei der
	Abfall in den unteren Schichten am stärksten ist.}
	\label{fig:barometric}
\end{figure}

\begin{figure}[h]
	\centering
	\includegraphics[width=0.85\textwidth]{plots/pressure_difference.pdf}
	\caption{Hydrostatische Druckdifferenz $\Delta p(\Delta h)$
	(Lemma~\ref{lem:hydrostatisch}, durchgezogen) im Vergleich zur
	Linearisierung $\dfrac{p_0}{H}\Delta h$
	(Lemma~\ref{lem:linearisierung}, gestrichelt) für
	$\Delta h \in [0, 3000]\ \mathrm{m}$. Für kleine $\Delta h$ stimmen
	beide Kurven nahezu überein; die Abweichung wächst quadratisch mit
	$\Delta h$, wie in Lemma~\ref{lem:linearisierung} formal gezeigt.}
	\label{fig:pressure_difference}
\end{figure}

\begin{figure}[h]
	\centering
	\includegraphics[width=0.85\textwidth]{plots/extraction_optimization.pdf}
	\caption{Normierte Leistungsfunktion $f(x) = x\sqrt{1-x}$
	(Satz~\ref{satz:optimal_x}). Die Funktion verschwindet an beiden
	Rändern ($f(0)=f(1)=0$) und besitzt ihr eindeutiges globales Maximum
	bei $x^\ast = 2/3$ mit $f(x^\ast) = 2\sqrt{3}/9 \approx 0{,}3849$.
	Die Kurve ist auf $(0, 2/3)$ streng monoton steigend und auf
	$(2/3, 1)$ streng monoton fallend, was die im Beweis verwendete
	Fallunterscheidung über das Vorzeichen von $f'$ visuell bestätigt.}
	\label{fig:optimization}
\end{figure}

\begin{figure}[h]
	\centering
	\includegraphics[width=0.85\textwidth]{plots/power_curves.pdf}
	\caption{Maximale Leistung $P_{\max}(\Delta p)$ nach
	Gleichung~\eqref{eq:pmax} für $\rho = 1{,}0\ \mathrm{kg/m^3}$ und
	Querschnittsflächen $A \in \{1, 5, 10, 20\}\ \mathrm{m^2}$. Die
	überproportionale ($\Delta p^{3/2}$-) Skalierung aus
	Satz~\ref{satz:skalierung} äußert sich in der zunehmenden
	Steigung der Kurven für größere $\Delta p$: Eine Verdoppelung von
	$\Delta p$ führt zu einer Leistungssteigerung um den Faktor
	$2^{3/2} \approx 2{,}83$, unabhängig von $A$.}
	\label{fig:power_curves}
\end{figure}

\begin{figure}[h]
	\centering
	\includegraphics[width=0.85\textwidth]{plots/equivalent_wind_speed.pdf}
	\caption{Äquivalente Windgeschwindigkeit $v_{\mathrm{eq}}(\Delta p) =
	\sqrt{2\Delta p/\rho}$ für $\rho = 1{,}0\ \mathrm{kg/m^3}$. Markiert
	ist der in den Beispielen~\ref{bsp:rotation},
	\ref{bsp:numerisch} und \ref{bsp:vergleich} verwendete Fall
	$\Delta p = 300\ \mathrm{Pa}$, der einer äquivalenten
	Windgeschwindigkeit von $v_{\mathrm{eq}} \approx 24{,}5\
	\mathrm{m/s}$ entspricht -- eine Größenordnung, die mit den in
	Abschnitt~\ref{sec:atmosphaere} diskutierten Aufwindgeschwindigkeiten
	in starken Cumulonimbus-Zellen konsistent ist.}
	\label{fig:equivalent_wind}
\end{figure}

\section{Zusammenfassung}

Diese Arbeit hat ein formales Modell für die Energiegewinnung aus
Druckdifferenzen in Wolken entwickelt. Auf Basis der Bernoulli-Gleichung
wurde gezeigt, dass:

\begin{itemize}
	\item die Leistung $P(x)$ eines Druckdifferenz-Wandlers durch die
	einparametrige Funktion $f(x) = x\sqrt{1-x}$ bestimmt wird, deren
	eindeutiges globales Maximum bei $x^\ast = 2/3$ liegt
	(Satz~\ref{satz:optimal_x});
	\item die maximale Leistung $P_{\max} = c_p A\sqrt{2/\rho}\,
	\Delta p^{3/2}$ mit $c_p = 2\sqrt{3}/9$ einer
	$\Delta p^{3/2}$-Skalierung folgt (Korollar~\ref{kor:pmax},
	Satz~\ref{satz:skalierung});
	\item hydrostatische Druckdifferenzen näherungsweise linear mit der
	Höhendifferenz $\Delta h$ wachsen, mit einer expliziten,
	quadratischen Fehlerschranke (Lemma~\ref{lem:linearisierung});
	\item die Leistungsdichte des Wandlers über den Koeffizienten $c_p$
	formal mit der kinetischen Energieflussdichte des äquivalenten
	Windes in Beziehung steht (Satz~\ref{satz:vergleich_wind}).
\end{itemize}

Damit erreicht das Modell eine geschlossene, vollständig formal
hergeleitete Beschreibung der Leistungsausbeute eines idealisierten
Druckdifferenz-Wandlers -- als Referenzpunkt für die Bewertung realer
technischer Umsetzungen.

\subsection{Anwendungen}

Das Modell eignet sich insbesondere für:

\begin{itemize}
	\item Theoretische Obergrenzen für geflesselte
	(\emph{tethered}) Energiewandler, die an Ballons oder Drachen in
	Wolkenhöhe betrieben werden;
	\item Vergleichende Bewertung verschiedener Standorte und
	Wolkentypen anhand der zu erwartenden Druckdifferenzen $\Delta p$;
	\item Dimensionierung von Wandler-Querschnittsflächen $A$ in
	Abhängigkeit von der Ziel-Leistung $P_{\max}$;
	\item Kombination mit bestehenden Airborne-Wind-Energy-Konzepten,
	bei denen sowohl kinetische als auch statische Druckanteile der
	Strömung genutzt werden können;
	\item Lehre und Veranschaulichung der Bernoulli-Gleichung und
	klassischer Optimierungsmethoden anhand eines konkreten,
	atmosphärenphysikalisch motivierten Beispiels.
\end{itemize}

\subsection{Ausblick}
\label{sec:ausblick}

Das in dieser Arbeit vorgestellte Modell ist ein idealisierter,
eindimensionaler Rahmen. Für eine technische Umsetzung sowie für die
wissenschaftliche Weiterentwicklung ergeben sich zahlreiche
Anschlussfragen, die im Folgenden ausführlich diskutiert werden.

\subsubsection{Technische Realisierung von Trägerplattformen}

Die Nutzung von Druckdifferenzen in Wolkenhöhe setzt eine Plattform
voraus, die den Wandler in der relevanten Höhe positioniert und dort
über längere Zeiträume stabil hält. Denkbare Ansätze umfassen:

\begin{itemize}
	\item \textbf{Geflesselte Ballonsysteme:} Ein mit Helium oder
	Heißluft gefüllter Auftriebskörper trägt den Wandler in
	Wolkenbasishöhe; die elektrische Leistung wird über ein
	leitfähiges Halteseil zum Boden übertragen. Die hydrostatische
	Komponente aus Abschnitt~\ref{sec:atmosphaere} (Größenordnung
	$10$--$30\ \mathrm{kPa}$ über $\Delta h \approx 1$--$3\ \mathrm{km}$)
	wäre hier potentiell relevant, sofern ein durchgängiger Kanal
	zwischen Boden- und Höhenniveau realisiert werden kann -- eine
	erhebliche konstruktive Herausforderung.
	\item \textbf{Fesseldrachen und Tragflächendrohnen:} Analog zu
	Crosswind-Kite-Systemen nach Loyd~\cite{loyd1980} könnten
	Tragflächen so positioniert werden, dass lokale dynamische
	Druckdifferenzen (Abschnitt~\ref{sec:atmosphaere},
	$\Delta p_{\mathrm{dyn}} \sim \tfrac{1}{2}\rho v^2$) über
	integrierte Wandlerkanäle abgegriffen werden.
	\item \textbf{Drohnenschwärme:} Mehrere autonome, höhenvariable
	Plattformen könnten räumlich verteilte Druckdifferenzzellen
	gleichzeitig vermessen und nutzen, wodurch sich ein Netzwerk von
	Wandlern ergibt, dessen Gesamtleistung sich (unter
	Vernachlässigung von Wechselwirkungen) additiv aus den einzelnen
	$P_{\max}$-Beiträgen ergibt.
\end{itemize}

\subsubsection{Materialwissenschaftliche und konstruktive Herausforderungen}

Der Betrieb in und um Wolken stellt erhebliche Anforderungen an Materialien
und Konstruktion:

\begin{itemize}
	\item \textbf{Vereisung:} Unterkühlte Wassertröpfchen in
	Cumulonimbus-Wolken führen zu Eisansatz an Rotorblättern und
	Kanalwänden, was die effektive Querschnittsfläche $A$ reduziert und
	die Strömung störungsanfälliger macht.
	\item \textbf{Blitzschutz:} Hochreichende Konvektionswolken sind mit
	signifikanter Blitzaktivität verbunden; leitfähige Halteseile und
	metallische Wandlerkomponenten müssen entsprechend dimensioniert und
	geerdet werden.
	\item \textbf{Materialermüdung:} Turbulente Anströmung führt zu
	hochfrequenten Lastwechseln, die bei der Auslegung von Rotorblättern,
	Lagern und Halteseilen berücksichtigt werden müssen.
	\item \textbf{Korrosion:} Die hohe Luftfeuchtigkeit innerhalb von
	Wolken begünstigt Korrosionsprozesse an metallischen Bauteilen,
	insbesondere an elektrischen Kontakten und Generatoren.
\end{itemize}

\subsubsection{Regelungstechnik und Maximum-Power-Point-Tracking}

Eine zentrale praktische Implikation von Satz~\ref{satz:optimal_x} ist,
dass der optimale Extraktionsanteil $x^\ast = 2/3$ \emph{unabhängig} von
$\Delta p$, $\rho$ und $A$ ist. Dies vereinfacht die Regelungsstrategie
gegenüber klassischen Maximum-Power-Point-Tracking-Verfahren (MPPT), wie
sie etwa bei Photovoltaikanlagen eingesetzt werden, erheblich: Anstatt den
Arbeitspunkt kontinuierlich an wechselnde Umgebungsbedingungen anzupassen,
genügt es im idealisierten Modell, das Verhältnis von abgebautem zu
verfügbarem Druck konstant bei $2/3$ zu halten. In der Praxis wird diese
Aussage durch Reibungseffekte, Kompressibilität und instationäre Strömung
modifiziert; eine vollständige regelungstechnische Behandlung müsste
daher
\begin{itemize}
	\item Echtzeit-Druckmessungen an mehreren Punkten der
	Druckdifferenzzelle,
	\item eine dynamische Schätzung von $\rho$ in Abhängigkeit von Höhe
	und Temperatur, sowie
	\item eine adaptive Nachführung des effektiven Strömungswiderstands
	des Wandlers (z.\,B.\ über variable Schaufelwinkel)
\end{itemize}
integrieren, um $x^\ast = 2/3$ auch unter realen, zeitveränderlichen
Bedingungen näherungsweise zu realisieren.

\subsubsection{Stochastische und turbulente Strömungsmodelle}

Das vorliegende Modell setzt eine \emph{stationäre}, \emph{laminare}
Strömung voraus. Reale atmosphärische Strömungen in und um konvektive
Zellen sind jedoch hochgradig turbulent und instationär. Eine
Weiterentwicklung des Modells könnte:

\begin{itemize}
	\item $\Delta p$ und $\rho$ als zeitabhängige stochastische Prozesse
	$\Delta p(t)$, $\rho(t)$ modellieren, etwa als
	Ornstein-Uhlenbeck-Prozesse mit empirisch bestimmten Parametern aus
	Radiosonden- oder Doppler-Radar-Daten~\cite{markowski2010};
	\item die Erwartungswertbildung $\mathbb{E}[P_{\max}(\Delta p(t))]$
	anstelle der deterministischen Größe $P_{\max}(\Delta p)$
	betrachten, wobei wegen der Konvexität von $\Delta p \mapsto
	\Delta p^{3/2}$ nach der Jensen-Ungleichung
	$\mathbb{E}[\Delta p(t)^{3/2}] \geq \mathbb{E}[\Delta p(t)]^{3/2}$
	gilt -- Schwankungen in $\Delta p$ erhöhen also tendenziell die
	mittlere Leistung gegenüber einer rein auf dem Mittelwert basierten
	Abschätzung;
	\item Reynoldszahl-abhängige Korrekturterme für die
	Bernoulli-Gleichung einführen, um Reibungsverluste in Abhängigkeit
	von der Strömungsgeschwindigkeit zu quantifizieren;
	\item dreidimensionale CFD-Simulationen (Computational Fluid
	Dynamics) durchführen, um die Wechselwirkung des Wandlers mit der
	umgebenden Wolkenströmung realistisch abzubilden.
\end{itemize}

\subsubsection{Kompressible Strömungseffekte bei großen Druckdifferenzen}

Für die in dieser Arbeit betrachteten Druckdifferenzen
($\Delta p \lesssim 1\ \mathrm{kPa}$, entsprechend $v_{\mathrm{eq}}
\lesssim 45\ \mathrm{m/s}$ bei $\rho \approx 1\ \mathrm{kg/m^3}$) ist die
Mach-Zahl $\mathrm{Ma} = v/c$ (mit Schallgeschwindigkeit $c \approx 340\
\mathrm{m/s}$) durchgehend $\mathrm{Ma} < 0{,}15$, sodass die
Inkompressibilitätsannahme der Bernoulli-Gleichung gerechtfertigt ist. Für
die in Abschnitt~\ref{sec:atmosphaere} diskutierten hydrostatischen
Druckdifferenzen über mehrere Kilometer Höhendifferenz ($\Delta p \sim
10$--$30\ \mathrm{kPa}$) wäre $v_{\mathrm{eq}}$ deutlich höher und eine
Erweiterung des Modells auf die kompressible Bernoulli-Gleichung (bzw.\
die Energiegleichung für ideale Gase mit Enthalpietermen) notwendig. Eine
solche Erweiterung würde insbesondere
\begin{itemize}
	\item die Temperaturänderung der Luft beim Druckabbau
	(adiabatische Expansion, $T \propto p^{(\gamma-1)/\gamma}$) und
	\item die damit verbundene Dichteänderung $\rho(p)$
\end{itemize}
berücksichtigen müssen, was die Leistungsfunktion $P(x)$ aus
Gleichung~\eqref{eq:leistung} nichtlinear in $x$ und $\Delta p$ koppelt
und voraussichtlich zu einem von $x^\ast = 2/3$ abweichenden optimalen
Extraktionsanteil führt.

\subsubsection{Skalierung zu Wandlerfeldern und Netzintegration}

Analog zu Windparks könnten mehrere Druckdifferenz-Wandler räumlich
verteilt betrieben werden ("Druckdifferenz-Farmen"). Hierbei stellen sich
Fragen der
\begin{itemize}
	\item gegenseitigen aerodynamischen Beeinflussung benachbarter
	Wandler (Nachlaufeffekte, analog zum Wake-Effekt bei Windturbinen);
	\item optimalen räumlichen Anordnung zur Maximierung der
	Gesamtleistung unter Berücksichtigung der räumlichen Korrelation von
	$\Delta p$-Feldern in konvektiven Zellen;
	\item Netzintegration der erzeugten Leistung, die aufgrund der
	Intermittenz konvektiver Ereignisse stark fluktuiert und daher
	Speicherlösungen (Batterien, Wasserstoff, Pumpspeicher) oder eine
	Kombination mit grundlastfähigen Quellen erfordert.
\end{itemize}

\subsubsection{Wechselwirkung mit der Wolkendynamik}

Eine grundsätzliche, bislang nicht quantifizierte Frage betrifft die
Rückwirkung der Energieentnahme auf die Wolkendynamik selbst. Die
Entnahme von Leistung $P_{\max}$ aus einer Druckdifferenzzelle entspricht
einer Reduktion der kinetischen bzw.\ potentiellen Energie der lokalen
Strömung. Bei hinreichend großer installierter Leistung relativ zur
Gesamtenergiebilanz der konvektiven Zelle könnte dies
\begin{itemize}
	\item die Aufwindgeschwindigkeit und damit $\Delta p_{\mathrm{dyn}}$
	selbst reduzieren (negative Rückkopplung),
	\item die vertikale Entwicklung der Wolke (Wolkenoberkante,
	Niederschlagsbildung) beeinflussen,
	\item im Extremfall mesoskalige Zirkulationsmuster verändern.
\end{itemize}
Eine quantitative Untersuchung dieser Rückkopplungseffekte würde
gekoppelte Modelle aus Wolkenphysik~\cite{houze2014,emanuel1994} und dem
hier entwickelten Wandlermodell erfordern, etwa in Form von
Sensitivitätsstudien mit numerischen Wettermodellen, die einen zusätzlichen
Energiesenkenterm proportional zu $P_{\max}(\Delta p(\mathbf{x},t))$
enthalten.

\subsubsection{Rechtliche und regulatorische Aspekte}

Der Betrieb geflesselter oder autonomer Plattformen in Höhen von mehreren
hundert bis mehreren tausend Metern unterliegt luftverkehrsrechtlichen
Bestimmungen (Luftraumklassifikation, Kennzeichnungspflichten,
Betriebsgenehmigungen). Eine praktische Umsetzung würde eine enge
Abstimmung mit Luftfahrtbehörden erfordern, insbesondere hinsichtlich:
\begin{itemize}
	\item der maximal zulässigen Betriebshöhe und der Abgrenzung zu
	kontrolliertem Luftraum,
	\item der Kennzeichnung und Beleuchtung der Plattformen zur
	Kollisionsvermeidung,
	\item der Wetterabhängigkeit des Betriebs (Einschränkungen bei
	Gewitterwarnungen, die naturgemäß genau die Bedingungen betreffen,
	unter denen $\Delta p$ am größten ist).
\end{itemize}

\subsubsection{Vergleich und Kombination mit Airborne Wind Energy}

Das Forschungsfeld der Airborne Wind Energy (AWE) befasst sich seit den
grundlegenden Arbeiten von Loyd~\cite{loyd1980} mit der Nutzung von
Höhenwinden über geflesselte Tragflächen und Drachen. Übersichtsarbeiten
wie Cherubini et al.~\cite{cherubini2015} und globale Potenzialanalysen
wie Archer und Caldeira~\cite{archer2009} zeigen, dass in Höhen von
$500$--$10\,000\ \mathrm{m}$ erhebliche kinetische Windenergiepotenziale
existieren. Das in dieser Arbeit vorgestellte Druckdifferenz-Modell
ergänzt diese Ansätze um eine \emph{statische} Komponente: Während AWE-
Systeme primär die kinetische Energie des Horizontalwinds nutzen, zielt
der Druckdifferenz-Wandler auf vertikale bzw.\ lokale Druckgradienten ab.
Eine Kombination beider Prinzipien -- etwa eine AWE-Plattform mit
integriertem Druckdifferenz-Kanal -- könnte die Gesamtenergieausbeute pro
Plattform erhöhen, wobei die in Satz~\ref{satz:vergleich_wind}
hergeleitete Beziehung $\pi_p = c_p \cdot \pi_w$ als Anhaltspunkt für die
relative Größenordnung beider Beiträge dienen kann.

\subsubsection{Experimentelle Validierung}

Vor einer technischen Umsetzung wäre eine experimentelle Validierung des
Modells in mehreren Stufen sinnvoll:
\begin{itemize}
	\item \textbf{Windkanalversuche:} Nachbildung der
	Druckdifferenzzelle in einem Windkanal mit kontrollierbarem
	$\Delta p$, um die Beziehung $P(x)$ aus Gleichung~\eqref{eq:leistung}
	und insbesondere das Maximum bei $x^\ast = 2/3$ experimentell zu
	überprüfen.
	\item \textbf{Feldmessungen mit Sondenballons:} Gleichzeitige Druck-
	und Windmessungen an mehreren Höhen innerhalb realer Cumulus- und
	Cumulonimbus-Wolken, um die in Abschnitt~\ref{sec:atmosphaere}
	angenommenen Größenordnungen für $\Delta p$ zu verifizieren.
	\item \textbf{Prototyp-Wandler:} Bau eines kleinskaligen
	Demonstrators (z.\,B.\ $A < 1\ \mathrm{m^2}$) zur Messung der
	tatsächlich erzielbaren Leistung im Vergleich zur theoretischen
	Obergrenze $P_{\max}$, woraus sich ein empirischer
	Gesamtwirkungsgrad $\eta_{\mathrm{real}} = P_{\mathrm{real}} /
	P_{\max}$ ableiten ließe.
\end{itemize}

\subsubsection{Weiterführende theoretische Fragestellungen}

Schließlich ergeben sich mehrere rein theoretische Anschlussfragen:

\begin{itemize}
	\item \textbf{Mehrpunkt-Wandlernetzwerke:} Statt einer einzelnen
	Druckdifferenzzelle könnte ein Netzwerk von Messpunkten betrachtet
	werden, das als gewichteter Graph modelliert wird, wobei Knoten
	Messpunkten und Kanten den paarweisen Druckdifferenzen entsprechen.
	Die Identifikation strukturell ähnlicher Teilnetzwerke -- etwa um
	wiederkehrende, energetisch günstige Konfigurationen über
	verschiedene Wolkensysteme hinweg zu erkennen -- könnte
	methodisch an den Subgraph Algorithmus anknüpfen, dessen
	Signatur-basierter Vergleich von Adjazenzstrukturen sich
	grundsätzlich auch auf gewichtete Druckdifferenzgraphen übertragen
	ließe.
	\item \textbf{Mehrdimensionale Optimierung:} Erweiterung der
	Optimierung von $x$ auf ein Paar $(x, \theta)$, wobei $\theta$ einen
	Anströmwinkel des Wandlers relativ zur lokalen Strömungsrichtung
	beschreibt, mit entsprechend erweiterter Bernoulli-Gleichung.
	\item \textbf{Zeitintegrierte Energiebilanzen:} Ãœbergang von der
	momentanen Leistung $P_{\max}(\Delta p)$ zu einer über Tages- oder
	Jahresgänge integrierten Energieausbeute $E = \int P_{\max}(\Delta
	p(t))\,dt$, unter Verwendung klimatologischer Statistiken zur
	Häufigkeitsverteilung von $\Delta p$ an einem gegebenen Standort.
	\item \textbf{Unsicherheitsquantifizierung:} Formale Fehlerschranken
	(analog zu Lemma~\ref{lem:linearisierung}) für $P_{\max}$ unter
	Messunsicherheiten in $\Delta p$, $\rho$ und $A$, um Konfidenzintervalle
	für prognostizierte Leistungswerte anzugeben.
\end{itemize}

Insgesamt zeigt diese Arbeit, dass bereits ein vergleichsweise einfaches,
vollständig formal beweisbares Modell -- basierend auf einer einzigen
Bernoulli-Gleichung und einer eindimensionalen Optimierung -- eine
geschlossene und überraschend reichhaltige Theorie der Energiegewinnung
aus Druckdifferenzen in Wolken liefert. Die zahlreichen oben skizzierten
Erweiterungsrichtungen zeigen zugleich, dass dieses einfache Modell
lediglich den Ausgangspunkt für eine deutlich umfassendere, interdisziplinäre
Forschungsagenda an der Schnittstelle von Strömungsmechanik,
Wolkenphysik, Regelungstechnik und erneuerbaren Energien darstellt.

\begin{thebibliography}{99}

\bibitem{holton2013}
Holton, J.\,R., Hakim, G.\,J.:
\emph{An Introduction to Dynamic Meteorology}, 5.\ Auflage.
Academic Press, 2013.

\bibitem{markowski2010}
Markowski, P., Richardson, Y.:
\emph{Mesoscale Meteorology in Midlatitudes}.
Wiley-Blackwell, 2010.

\bibitem{houze2014}
Houze, R.\,A.:
\emph{Cloud Dynamics}, 2.\ Auflage.
Academic Press, 2014.

\bibitem{emanuel1994}
Emanuel, K.\,A.:
\emph{Atmospheric Convection}.
Oxford University Press, 1994.

\bibitem{stull1988}
Stull, R.\,B.:
\emph{An Introduction to Boundary Layer Meteorology}.
Kluwer Academic Publishers, 1988.

\bibitem{betz1920}
Betz, A.:
Das Maximum der theoretisch möglichen Ausnutzung des Windes durch
Windmotoren.
\emph{Zeitschrift für das gesamte Turbinenwesen}, 26, 1920.

\bibitem{loyd1980}
Loyd, M.\,L.:
Crosswind kite power.
\emph{Journal of Energy}, 4(3), 1980.

\bibitem{archer2009}
Archer, C.\,L., Caldeira, K.:
Global Assessment of High-Altitude Wind Power.
\emph{Energies}, 2(2), 2009.

\bibitem{cherubini2015}
Cherubini, A., Papini, A., Vertechy, R., Fontana, M.:
Airborne Wind Energy Systems: A review of the technologies.
\emph{Renewable and Sustainable Energy Reviews}, 51, 2015.

\bibitem{epp2026subgraph}
Epp, S.:
\emph{Der Subgraph Algorithmus}.
Eigene Referenzarbeit zur Methodik formaler Optimalitätsbeweise, 2026.

\end{thebibliography}

\end{document}

